\documentclass[10pt,oneside,openany]{memoir}
%%%%%%%%%%%%%%%%%%%%%%%%%%%%%%%%%%%%%%%%%%%%%%%%%%%%%%%%%%%%%%%%%%%%%%%%%%%%%%%%%%%%%
% Font options
%%%%%%%%%%%%%%%%%%%%%%%%%%%%%%%%%%%%%%%%%%%%%%%%%%%%%%%%%%%%%%%%%%%%%%%%%%%%%%%%%%%%%

\iffalse
%font style 1
\usepackage[p,osf]{scholax}
\usepackage{amsmath,amsthm,amssymb}
\usepackage[scaled=1.075,ncf,vvarbb]{newtxmath}
\fi

%font style2
\usepackage[T1]{fontenc}
\usepackage{amsmath}
\usepackage[fixamsmath]{mathtools}  % Extension to amsmath
\usepackage{amsthm}
\usepackage{amssymb}
\renewcommand*{\mathbf}[1]{\mathbb{#1}}
\renewcommand{\qedsymbol}{$\blacksquare$}

%font style 3
%\usepackage{newpxtext}
%\usepackage{eulerpx}
%\usepackage{mathrsfs}
%\usepackage{eucal}

\usepackage[uprightgreeks,light]{kpfonts}
\usepackage{datetime}
%%%%%%%%%%%%%%%%%%%%%%%%%%%%%%%%%%%%%%%%%%%%%%%%%%%%%%%%%%%%%%%%%%%%%%%%%%%%%%%%%%%%%
% Page layout and headers
%%%%%%%%%%%%%%%%%%%%%%%%%%%%%%%%%%%%%%%%%%%%%%%%%%%%%%%%%%%%%%%%%%%%%%%%%%%%%%%%%%%%%


\usepackage[margin=1in]{geometry}
\usepackage{graphicx} % needed for \scalebox
\usepackage{fancyhdr}
\usepackage{datetime}

\pagestyle{fancy}
\fancyhf{}
\cfoot{\scriptsize\thepage}
\fancyhead[R]{\scalebox{0.8}{\nouppercase{\rightmark}}}
\fancyhead[L]{\scalebox{0.8}{\nouppercase{\leftmark}}}

\fancypagestyle{plain}{%
    \fancyhf{}
    \cfoot{\scriptsize\thepage}
    \fancyhead[R]{\scalebox{0.8}{\phantom{a}}}
    \fancyhead[L]{\scalebox{0.8}{Applied Analysis}}
}

\renewcommand{\headrulewidth}{0pt}
\renewcommand{\footrulewidth}{0pt}

%%%%%%%%%%%%%%%%%%%%%%%%%%%%%%%%%%%%%%%%%%%%%%%%%%%%%%%%%%%%%%%%%%%%%%%%%%%%%%%%%%%%%
% Theorem environments
%%%%%%%%%%%%%%%%%%%%%%%%%%%%%%%%%%%%%%%%%%%%%%%%%%%%%%%%%%%%%%%%%%%%%%%%%%%%%%%%%%%%%

\usepackage{keytheorems}

\newkeytheoremstyle{def}{
    inherit-style=definition,
    spaceabove=1.5em,
    spacebelow=1.5em,
    postheadspace=0.5em,
}

\newkeytheoremstyle{thm}{
    spaceabove=1.5em,
    spacebelow=1.5em,
    postheadspace=0.5em,
}

\newkeytheorem{theorem}[
    name = Theorem,
    style=thm,
    numberwithin = section,
]

\newkeytheorem{theorem*}[
    name=Theorem,
    style=thm,
    numbered=false
]

\newkeytheorem{proposition}[
    name = Proposition,
    style=thm,
    numberwithin = section,
    sibling=theorem
]

\newkeytheorem{lemma}[
    name = Lemma,
    style=thm,
    numberwithin = section,
    sibling=theorem
]

\newkeytheorem{corollary}[
    name = Corollary,
    style=thm,
    numberwithin = section,
    sibling=theorem
]

\newkeytheorem{definition}[
    name = Definition,
    style=def,
    numberwithin = section,
    style = definition
]

\newkeytheorem{example}[
    name = Example,
    style=def,
    numberwithin = section,
    style = definition
]

\newkeytheorem{remark}[
    name=Remark,
    style=remark,
    numbered=false
]

\newkeytheorem{question}[
    name=Question,
    style=remark,
    numbered=false
]

\newkeytheorem{exercise}[
    name = Exercise,
    style=thm
]
\newkeytheorem{problem}[
    name = Problem,
    style=thm
]

\newkeytheorem{axiom}[name = Axiom]

%%%%%%%%%%%%%%%%%%%%%%%%%%%%%%%%%%%%%%%%%%%%%%%%%%%%%%%%%%%%%%%%%%%%%%%%%%%%%%%%%%%%%
% Section and chapter formatting
%%%%%%%%%%%%%%%%%%%%%%%%%%%%%%%%%%%%%%%%%%%%%%%%%%%%%%%%%%%%%%%%%%%%%%%%%%%%%%%%%%%%%

\usepackage{enumitem}
\usepackage[explicit]{titlesec}

\titleformat{\chapter}[display]
    {\bfseries\LARGE\raggedright}
    {Chapter {\thechapter}}
    {.1ex minus .1ex}
    {\Large #1}

\titlespacing{\chapter}
    {3pc}{*3}{40pt}[3pc]

\titleformat{\section}[block]
    {\normalfont\bfseries\large}
    {}{0pt}
    {\fbox{\strut \S\ \thesection.\ #1}}

\titlespacing{\section}
    {0pt}{3ex plus .1ex minus .2ex}{3ex plus .1ex minus .2ex}

\setlength{\fboxsep}{4pt}    % padding inside the box
\setlength{\fboxrule}{0.6pt} % border thickness

\titleformat{name=\section,numberless}[block]
    {\normalfont\bfseries\large\raggedright}
    {}{0pt}
    {\fbox{\strut #1}}

%%%%%%%%%%%%%%%%%%%%%%%%%%%%%%%%%%%%%%%%%%%%%%%%%%%%%%%%%%%%%%%%%%%%%%%%%%%%%%%%%%%%%
% Graphics, symbols, and miscellaneous packages
%%%%%%%%%%%%%%%%%%%%%%%%%%%%%%%%%%%%%%%%%%%%%%%%%%%%%%%%%%%%%%%%%%%%%%%%%%%%%%%%%%%%%

\usepackage{tikz}
\usepackage{tikz-cd}
\usepackage{float}
\usepackage{wasysym}
\usepackage{hyperref}
\usepackage{ulem}
\usepackage{pgfplots}
\usetikzlibrary{patterns}
\usetikzlibrary{positioning,arrows.meta}
\usetikzlibrary{decorations.pathreplacing}
\usetikzlibrary{patterns}
\usepgfplotslibrary{fillbetween}
\pgfplotsset{compat=1.18}
\usepackage{etoolbox}
\usepackage{bm}

%%%%%%%%%%%%%%%%%%%%%%%%%%%%%%%%%%%%%%%%%%%%%%%%%%%%%%%%%%%%%%%%%%%%%%%%%%%%%%%%%%%%%
% Final Formatting/Extra commands
%%%%%%%%%%%%%%%%%%%%%%%%%%%%%%%%%%%%%%%%%%%%%%%%%%%%%%%%%%%%%%%%%%%%%%%%%%%%%%%%%%%%%

\newcommand{\homeworktitle}[3]{%
    \begin{center}
        {\large #1 \\[0.1in] #2}\\[.175in]
        {Name:} {\underline{#3\hspace*{2in}}}\\[0.15in]
    \end{center}
    \vspace{4pt}
}
\newdateformat{specialdate}{\THEYEAR\ \monthname\ \THEDAY}

%%%%%%%%%%%%%%%%%%%%%%%%%%%%%%%%%%%%%%%%%%%%%%%%%%%%%%%%%%%%%
%%%%%%%%%%%%%%%%%%%%%%%%%%%%%%%%%%%%%%%%%%%%%%%%%%%%%%%%%%%%%

% ============================================================
% makros.tex — reorganized, full restored version
% ============================================================

% ============================================================
% Sha / Cyrillic support notes
% ============================================================

%to make the correct symbol for Sha
%\newcommand\cyr{%
%\renewcommand\rmdefault{wncyr}%
%\renewcommand\sfdefault{wncyss}%
%\renewcommand\encodingdefault{OT2}%
%\normalfont \selectfont} \DeclareTextFontCommand{\textcyr}{\cyr}

% ============================================================
% Math operators
% ============================================================

\DeclareMathOperator{\ab}{ab}
\DeclareMathOperator{\ad}{ad}
\DeclareMathOperator{\adj}{adj}
\DeclareMathOperator{\alg}{alg}
\DeclareMathOperator{\Alt}{Alt}
\DeclareMathOperator{\Ann}{Ann}
\DeclareMathOperator{\arith}{arith}
\DeclareMathOperator{\Aut}{Aut}
\DeclareMathOperator{\Be}{B}
\DeclareMathOperator{\Bd}{Bd}
\DeclareMathOperator{\card}{card}
\DeclareMathOperator{\Char}{char}
\DeclareMathOperator{\csp}{csp}
\DeclareMathOperator{\codim}{codim}
\DeclareMathOperator{\coker}{coker}
\DeclareMathOperator{\coh}{H}
\DeclareMathOperator{\compl}{compl}
\DeclareMathOperator{\conj}{conj}
\DeclareMathOperator{\cont}{cont}
\DeclareMathOperator{\Cov}{Cov}
\DeclareMathOperator{\crys}{crys}
\DeclareMathOperator{\Crys}{Crys}
\DeclareMathOperator{\cusp}{cusp}
\DeclareMathOperator{\diag}{diag}
\DeclareMathOperator{\diam}{diam}
\DeclareMathOperator{\Dom}{Dom}
\DeclareMathOperator{\disc}{disc}
\DeclareMathOperator{\dist}{dist}
\DeclareMathOperator{\Div}{div}
\DeclareMathOperator{\dR}{dR}
\DeclareMathOperator{\Eis}{Eis}
\DeclareMathOperator{\End}{End}
\DeclareMathOperator{\ev}{ev}
\DeclareMathOperator{\eval}{eval}
\DeclareMathOperator{\Eq}{Eq}
\DeclareMathOperator{\Ext}{Ext}
\DeclareMathOperator{\Fil}{Fil}
\DeclareMathOperator{\Fitt}{Fitt}
\DeclareMathOperator{\Frob}{Frob}
\DeclareMathOperator{\G}{G}
\DeclareMathOperator{\Gal}{Gal}
\DeclareMathOperator{\GL}{GL}
\DeclareMathOperator{\Gr}{Gr}
\DeclareMathOperator{\Graph}{Graph}
\DeclareMathOperator{\GSp}{GSp}
\DeclareMathOperator{\GUn}{GU}
\DeclareMathOperator{\Hom}{Hom}
\DeclareMathOperator{\id}{id}
\DeclareMathOperator{\Id}{Id}
\DeclareMathOperator{\Ik}{Ik}
\DeclareMathOperator{\IM}{Im}
\DeclareMathOperator{\Image}{im}
\DeclareMathOperator{\Ind}{Ind}
\DeclareMathOperator{\Inf}{inf}
\DeclareMathOperator{\ins}{ins}
\DeclareMathOperator{\inv}{inv}
\DeclareMathOperator{\Isom}{Isom}
\DeclareMathOperator{\J}{J}
\DeclareMathOperator{\Jac}{Jac}
\DeclareMathOperator{\lcm}{lcm}
\DeclareMathOperator{\length}{length}
\DeclareMathOperator*{\limit}{limit}
\DeclareMathOperator{\Log}{Log}
\DeclareMathOperator{\M}{M}
\DeclareMathOperator{\Mat}{Mat}
\DeclareMathOperator{\N}{N}
\DeclareMathOperator{\Nm}{Nm}
\DeclareMathOperator{\NIk}{N-Ik}
\DeclareMathOperator{\NSK}{N-SK}
\DeclareMathOperator{\new}{new}
\DeclareMathOperator{\obj}{obj}
\DeclareMathOperator{\old}{old}
\DeclareMathOperator{\ord}{ord}
\DeclareMathOperator{\Or}{O}
\DeclareMathOperator{\op}{op}
\DeclareMathOperator{\out}{out}
\DeclareMathOperator{\PGL}{PGL}
\DeclareMathOperator{\PGSp}{PGSp}
\DeclareMathOperator{\rank}{rank}
\DeclareMathOperator{\Ran}{Ran}
\DeclareMathOperator{\rel}{Rel}
\DeclareMathOperator{\Real}{Re}
\DeclareMathOperator{\RES}{res}
\DeclareMathOperator{\Res}{Res}
%\DeclareMathOperator{\Sha}{\textcyr{Sh}}
\DeclareMathOperator{\Sel}{Sel}
\DeclareMathOperator{\semi}{ss}
\DeclareMathOperator{\sgn}{sign}
\DeclareMathOperator{\SK}{SK}
\DeclareMathOperator{\SL}{SL}
\DeclareMathOperator{\SO}{SO}
\DeclareMathOperator{\Sp}{Sp}
\DeclareMathOperator{\Span}{span}
\DeclareMathOperator{\Spec}{Spec}
\DeclareMathOperator{\spin}{spin}
\DeclareMathOperator{\st}{st}
\DeclareMathOperator{\St}{St}
\DeclareMathOperator{\SUn}{SU}
\DeclareMathOperator{\supp}{supp}
\DeclareMathOperator{\Sup}{sup}
\DeclareMathOperator{\Sym}{Sym}
\DeclareMathOperator{\Tam}{Tam}
\DeclareMathOperator{\tors}{tors}
\DeclareMathOperator{\tr}{tr}
\DeclareMathOperator{\Tr}{Tr}
\DeclareMathOperator{\un}{un}
\DeclareMathOperator{\Un}{U}
\DeclareMathOperator{\val}{val}
\DeclareMathOperator{\vol}{vol}

% ============================================================
% General aliases and short macros
% ============================================================

\newcommand{\absgal}{\G_{\bbQ}}
\newcommand{\Loc}{L_{\operatorname{loc}}}
\renewcommand{\k}{\kappa}
\newcommand{\Ff}{F_{f}}
%\newcommand{\ts}{\,^{t}\!}
\newcommand{\lat}{\mathscr{L}}
\newcommand{\ov}[1]{\overline{#1}}
\newcommand{\up}{\upsilon}
\newcommand{\e}{\epsilon}
\newcommand{\tac}{\textasteriskcentered}

%mahesh macros
\newcommand{\tm}{\textrm}

\newcommand{\bone}{\mathbf{1}}
\newcommand{\sign}{\mathrm{sign}}
\newcommand{\eps}{\varepsilon}
\newcommand{\textui}[1]{\underline{\textit{#1}}}

%\newcommand{\ov}{\overline}
%\newcommand{\un}{\underline}
\newcommand{\fin}{\mathrm{fin}}
\newcommand{\nl}{\newline \mbox{}}
\newcommand{\h}[1]{\hspace{#1pt}}
\newcommand{\mtext}[1]{\hspace{6pt}\text{#1}\hspace{6pt}}
\newcommand{\lnorm}{\left\lVert}
\newcommand{\rnorm}{\right\rVert}
\newcommand{\ds}{\displaystyle}
\newcommand{\ts}{\textstyle}
\newcommand{\sfrac}[2]{{}^{#1}\mskip -5mu/\mskip -3mu_{#2}}

% ============================================================
% Category-style operators and contoured variants
% ============================================================

\DeclareMathOperator{\Sets}{S \mkern1.04mu e \mkern1.04mu t \mkern1.04mu s}
    \newcommand{\cSets}{\scalebox{1.02}{\contour{black}{$\Sets$}}}
    
\DeclareMathOperator{\Groups}{G \mkern1.04mu r \mkern1.04mu o \mkern1.04mu u \mkern1.04mu p \mkern1.04mu s}
    \newcommand{\cGroups}{\scalebox{1.02}{\contour{black}{$\Groups$}}}

\DeclareMathOperator{\TTop}{T \mkern1.04mu o \mkern1.04mu p}
    \newcommand{\cTop}{\scalebox{1.02}{\contour{black}{$\TTop$}}}

\DeclareMathOperator{\Htp}{H \mkern1.04mu t \mkern1.04mu p}
    \newcommand{\cHtp}{\scalebox{1.02}{\contour{black}{$\Htp$}}}

\DeclareMathOperator{\Mod}{M \mkern1.04mu o \mkern1.04mu d}
    \newcommand{\cMod}{\scalebox{1.02}{\contour{black}{$\Mod$}}}

\DeclareMathOperator{\Ab}{A \mkern1.04mu b}
    \newcommand{\cAb}{\scalebox{1.02}{\contour{black}{$\Ab$}}}

\DeclareMathOperator{\Rings}{R \mkern1.04mu i \mkern1.04mu n \mkern1.04mu g \mkern1.04mu s}
    \newcommand{\cRings}{\scalebox{1.02}{\contour{black}{$\Rings$}}}

\DeclareMathOperator{\ComRings}{C \mkern1.04mu o \mkern1.04mu m \mkern1.04mu R \mkern1.04mu i \mkern1.04mu n \mkern1.04mu g \mkern1.04mu s}
    \newcommand{\cComRings}{\scalebox{1.05}{\contour{black}{$\ComRings$}}}

\DeclareMathOperator{\hHom}{H \mkern1.04mu o \mkern1.04mu m}
    \newcommand{\cHom}{\scalebox{1.02}{\contour{black}{$\hHom$}}}

% ============================================================
% Mathcal
% ============================================================

\newcommand{\cA}{\mathcal{A}}
\newcommand{\cB}{\mathcal{B}}
\newcommand{\cC}{\mathcal{C}}
\newcommand{\cD}{\mathcal{D}}
\newcommand{\cE}{\mathcal{E}}
\newcommand{\cF}{\mathcal{F}}
\newcommand{\cG}{\mathcal{G}}
\newcommand{\cH}{\mathcal{H}}
\newcommand{\cI}{\mathcal{I}}
\newcommand{\cJ}{\mathcal{J}}
\newcommand{\cK}{\mathcal{K}}
\newcommand{\cL}{\mathcal{L}}
\newcommand{\cM}{\mathcal{M}}
\newcommand{\cN}{\mathcal{N}}
\newcommand{\cO}{\mathcal{O}}
\newcommand{\cP}{\mathcal{P}}
\newcommand{\cQ}{\mathcal{Q}}
\newcommand{\cR}{\mathcal{R}}
\newcommand{\cS}{\mathcal{S}}
\newcommand{\cT}{\mathcal{T}}
\newcommand{\cU}{\mathcal{U}}
\newcommand{\cV}{\mathcal{V}}
\newcommand{\cW}{\mathcal{W}}
\newcommand{\cX}{\mathcal{X}}
\newcommand{\cY}{\mathcal{Y}}
\newcommand{\cZ}{\mathcal{Z}}

% ============================================================
% Mathscrl% 
%============================================================

\newcommand{\scra}{\mathscr{a}}
\newcommand{\scrA}{\mathscr{A}}
\newcommand{\scrb}{\mathscr{b}}
\newcommand{\scrB}{\mathscr{B}}
\newcommand{\scrc}{\mathscr{c}}
\newcommand{\scrC}{\mathscr{C}}
\newcommand{\scrd}{\mathscr{d}}
\newcommand{\scrD}{\mathscr{D}}
\newcommand{\scre}{\mathscr{e}}
\newcommand{\scrE}{\mathscr{E}}
\newcommand{\scrf}{\mathscr{f}}
\newcommand{\scrF}{\mathscr{F}}
\newcommand{\scrg}{\mathscr{g}}
\newcommand{\scrG}{\mathscr{G}}
\newcommand{\scrh}{\mathscr{h}}
\newcommand{\scrH}{\mathscr{H}}
\newcommand{\scri}{\mathscr{i}}
\newcommand{\scrI}{\mathscr{I}}
\newcommand{\scrj}{\mathscr{j}}
\newcommand{\scrJ}{\mathscr{J}}
\newcommand{\scrk}{\mathscr{k}}
\newcommand{\scrK}{\mathscr{K}}
\newcommand{\scrl}{\mathscr{l}}
\newcommand{\scrL}{\mathscr{L}}
\newcommand{\scrm}{\mathscr{m}}
\newcommand{\scrM}{\mathscr{M}}
\newcommand{\scrn}{\mathscr{n}}
\newcommand{\scrN}{\mathscr{N}}
\newcommand{\scro}{\mathscr{o}}
\newcommand{\scrO}{\mathscr{O}}
\newcommand{\scrp}{\mathscr{p}}
\newcommand{\scrP}{\mathscr{P}}
\newcommand{\scrq}{\mathscr{q}}
\newcommand{\scrQ}{\mathscr{Q}}
\newcommand{\scrr}{\mathscr{r}}
\newcommand{\scrR}{\mathscr{R}}
\newcommand{\scrs}{\mathscr{s}}
\newcommand{\scrS}{\mathscr{S}}
\newcommand{\scrt}{\mathscr{t}}
\newcommand{\scrT}{\mathscr{T}}
\newcommand{\scru}{\mathscr{u}}
\newcommand{\scrU}{\mathscr{U}}
\newcommand{\scrv}{\mathscr{v}}
\newcommand{\scrV}{\mathscr{V}}
\newcommand{\scrw}{\mathscr{w}}
\newcommand{\scrW}{\mathscr{W}}
\newcommand{\scrx}{\mathscr{x}}
\newcommand{\scrX}{\mathscr{X}}
\newcommand{\scry}{\mathscr{y}}
\newcommand{\scrY}{\mathscr{Y}}
\newcommand{\scrz}{\mathscr{z}}
\newcommand{\scrZ}{\mathscr{Z}}

% ============================================================
% Mathfrak (missing \fi)
% ============================================================

\newcommand{\fa}{\mathfrak{a}}
\newcommand{\fA}{\mathfrak{A}}
\newcommand{\fb}{\mathfrak{b}}
\newcommand{\fB}{\mathfrak{B}}
\newcommand{\fc}{\mathfrak{c}}
\newcommand{\fC}{\mathfrak{C}}
\newcommand{\fd}{\mathfrak{d}}
\newcommand{\fD}{\mathfrak{D}}
\newcommand{\fe}{\mathfrak{e}}
\newcommand{\fE}{\mathfrak{E}}
\newcommand{\ff}{\mathfrak{f}}
\newcommand{\fF}{\mathfrak{F}}
\newcommand{\fg}{\mathfrak{g}}
\newcommand{\fG}{\mathfrak{G}}
\newcommand{\fh}{\mathfrak{h}}
\newcommand{\fH}{\mathfrak{H}}
\newcommand{\fI}{\mathfrak{I}}
\newcommand{\fj}{\mathfrak{j}}
\newcommand{\fJ}{\mathfrak{J}}
\newcommand{\fk}{\mathfrak{k}}
\newcommand{\fK}{\mathfrak{K}}
\newcommand{\fl}{\mathfrak{l}}
\newcommand{\fL}{\mathfrak{L}}
\newcommand{\fm}{\mathfrak{m}}
\newcommand{\fM}{\mathfrak{M}}
\newcommand{\fn}{\mathfrak{n}}
\newcommand{\fN}{\mathfrak{N}}
\newcommand{\fo}{\mathfrak{o}}
\newcommand{\fO}{\mathfrak{O}}
\newcommand{\fp}{\mathfrak{p}}
\newcommand{\fP}{\mathfrak{P}}
\newcommand{\fq}{\mathfrak{q}}
\newcommand{\fQ}{\mathfrak{Q}}
\newcommand{\fr}{\mathfrak{r}}
\newcommand{\fR}{\mathfrak{R}}
\newcommand{\fs}{\mathfrak{s}}
\newcommand{\fS}{\mathfrak{S}}
\newcommand{\ft}{\mathfrak{t}}
\newcommand{\fT}{\mathfrak{T}}
\newcommand{\fu}{\mathfrak{u}}
\newcommand{\fU}{\mathfrak{U}}
\newcommand{\fv}{\mathfrak{v}}
\newcommand{\fV}{\mathfrak{V}}
\newcommand{\fw}{\mathfrak{w}}
\newcommand{\fW}{\mathfrak{W}}
\newcommand{\fx}{\mathfrak{x}}
\newcommand{\fX}{\mathfrak{X}}
\newcommand{\fy}{\mathfrak{y}}
\newcommand{\fY}{\mathfrak{Y}}
\newcommand{\fz}{\mathfrak{z}}
\newcommand{\fZ}{\mathfrak{Z}}

% ============================================================
% Mathbf
% ============================================================

\newcommand{\bfA}{\mathbf{A}}
\newcommand{\bfB}{\mathbf{B}}
\newcommand{\bfC}{\mathbf{C}}
\newcommand{\bfD}{\mathbf{D}}
\newcommand{\bfE}{\mathbf{E}}
\newcommand{\bfF}{\mathbf{F}}
\newcommand{\bfG}{\mathbf{G}}
\newcommand{\bfH}{\mathbf{H}}
\newcommand{\bfI}{\mathbf{I}}
\newcommand{\bfJ}{\mathbf{J}}
\newcommand{\bfK}{\mathbf{K}}
\newcommand{\bfL}{\mathbf{L}}
\newcommand{\bfM}{\mathbf{M}}
\newcommand{\bfN}{\mathbf{N}}
\newcommand{\bfO}{\mathbf{O}}
\newcommand{\bfP}{\mathbf{P}}
\newcommand{\bfQ}{\mathbf{Q}}
\newcommand{\bfR}{\mathbf{R}}
\newcommand{\bfS}{\mathbf{S}}
\newcommand{\bfT}{\mathbf{T}}
\newcommand{\bfU}{\mathbf{U}}
\newcommand{\bfV}{\mathbf{V}}
\newcommand{\bfW}{\mathbf{W}}
\newcommand{\bfX}{\mathbf{X}}
\newcommand{\bfY}{\mathbf{Y}}
\newcommand{\bfZ}{\mathbf{Z}}

\newcommand{\bfa}{\mathbf{a}}
\newcommand{\bfb}{\mathbf{b}}
\newcommand{\bfc}{\mathbf{c}}
\newcommand{\bfd}{\mathbf{d}}
\newcommand{\bfe}{\mathbf{e}}
\newcommand{\bff}{\mathbf{f}}
\newcommand{\bfg}{\mathbf{g}}
\newcommand{\bfh}{\mathbf{h}}
\newcommand{\bfi}{\mathbf{i}}
\newcommand{\bfj}{\mathbf{j}}
\newcommand{\bfk}{\mathbf{k}}
\newcommand{\bfl}{\mathbf{l}}
\newcommand{\bfm}{\mathbf{m}}
\newcommand{\bfn}{\mathbf{n}}
\newcommand{\bfo}{\mathbf{o}}
\newcommand{\bfp}{\mathbf{p}}
\newcommand{\bfq}{\mathbf{q}}
\newcommand{\bfr}{\mathbf{r}}
\newcommand{\bfs}{\mathbf{s}}
\newcommand{\bft}{\mathbf{t}}
\newcommand{\bfu}{\mathbf{u}}
\newcommand{\bfv}{\mathbf{v}}
\newcommand{\bfw}{\mathbf{w}}
\newcommand{\bfx}{\mathbf{x}}
\newcommand{\bfy}{\mathbf{y}}
\newcommand{\bfz}{\mathbf{z}}

% ============================================================
% Blackboard bold
% ============================================================

\newcommand{\bbA}{\mathbb{A}}
\newcommand{\bbB}{\mathbb{B}}
\newcommand{\bbC}{\mathbb{C}}
\newcommand{\bbD}{\mathbb{D}}
\newcommand{\bbE}{\mathbb{E}}
\newcommand{\bbF}{\mathbb{F}}
\newcommand{\bbG}{\mathbb{G}}
\newcommand{\bbH}{\mathbb{H}}
\newcommand{\bbI}{\mathbb{I}}
\newcommand{\bbJ}{\mathbb{J}}
\newcommand{\bbK}{\mathbb{K}}
\newcommand{\bbL}{\mathbb{L}}
\newcommand{\bbM}{\mathbb{M}}
\newcommand{\bbN}{\mathbb{N}}
\newcommand{\bbO}{\mathbb{O}}
\newcommand{\bbP}{\mathbb{P}}
\newcommand{\bbQ}{\mathbb{Q}}
\newcommand{\bbR}{\mathbb{R}}
\newcommand{\bbS}{\mathbb{S}}
\newcommand{\bbT}{\mathbb{T}}
\newcommand{\bbU}{\mathbb{U}}
\newcommand{\bbV}{\mathbb{V}}
\newcommand{\bbW}{\mathbb{W}}
\newcommand{\bbX}{\mathbb{X}}
\newcommand{\bbY}{\mathbb{Y}}
\newcommand{\bbZ}{\mathbb{Z}}

\newcommand{\Bmu}{\mbox{$\raisebox{-0.59ex}
  {$l$}\hspace{-0.18em}\mu\hspace{-0.88em}\raisebox{-0.98ex}{\scalebox{2}
  {$\color{white}.$}}\hspace{-0.416em}\raisebox{+0.88ex}
  {$\color{white}.$}\hspace{0.46em}$}{}}  %need graphicx and xcolor. this produces blackboard bold mu 

% ============================================================
% Greek-letter aliases
% ============================================================

\newcommand{\jota}{\jmath}

% ============================================================
% Matrices
% ============================================================

\newcommand{\bmat}{\left( \begin{matrix}}
\newcommand{\emat}{\end{matrix} \right)}

\newcommand{\bbmat}{\left[ \begin{matrix}}
\newcommand{\ebmat}{\end{matrix} \right]}

\newcommand{\pmat}{\left( \begin{smallmatrix}}
\newcommand{\epmat}{\end{smallmatrix} \right)}

\newcommand{\mat}[4]{\begin{pmatrix}{#1}&{#2}\\{#3}&{#4}\end{pmatrix}}

% ============================================================
% Residues and averaged integrals
% ============================================================

\newcommand{\res}[1]{\underset{#1}{\RES}\,}

\makeatletter
\newcommand*{\mint}[1]{%
  % #1: overlay symbol
  \mint@l{#1}{}%
}
\newcommand*{\mint@l}[2]{%
  % #1: overlay symbol
  % #2: limits
  \@ifnextchar\limits{%
    \mint@l{#1}%
  }{%
    \@ifnextchar\nolimits{%
      \mint@l{#1}%
    }{%
      \@ifnextchar\displaylimits{%
        \mint@l{#1}%
      }{%
        \mint@s{#2}{#1}%
      }%
    }%
  }%
}
\newcommand*{\mint@s}[2]{%
  % #1: limits
  % #2: overlay symbol
  \@ifnextchar_{%
    \mint@sub{#1}{#2}%
  }{%
    \@ifnextchar^{%
      \mint@sup{#1}{#2}%
    }{%
      \mint@{#1}{#2}{}{}%
    }%
  }%
}
\def\mint@sub#1#2_#3{%
  \@ifnextchar^{%
    \mint@sub@sup{#1}{#2}{#3}%
  }{%
    \mint@{#1}{#2}{#3}{}%
  }%
}
\def\mint@sup#1#2^#3{%
  \@ifnextchar_{%
    \mint@sup@sub{#1}{#2}{#3}%
  }{%
    \mint@{#1}{#2}{}{#3}%
  }%
}
\def\mint@sub@sup#1#2#3^#4{%
  \mint@{#1}{#2}{#3}{#4}%
}
\def\mint@sup@sub#1#2#3_#4{%
  \mint@{#1}{#2}{#4}{#3}%
}
\newcommand*{\mint@}[4]{%
  % #1: \limits, \nolimits, \displaylimits
  % #2: overlay symbol: -, =, ...
  % #3: subscript
  % #4: superscript
  \mathop{}%
  \mkern-\thinmuskip
  \mathchoice{%
    \mint@@{#1}{#2}{#3}{#4}%
        \displaystyle\textstyle\scriptstyle
  }{%
    \mint@@{#1}{#2}{#3}{#4}%
        \textstyle\scriptstyle\scriptstyle
  }{%
    \mint@@{#1}{#2}{#3}{#4}%
        \scriptstyle\scriptscriptstyle\scriptscriptstyle
  }{%
    \mint@@{#1}{#2}{#3}{#4}%
        \scriptscriptstyle\scriptscriptstyle\scriptscriptstyle
  }%
  \mkern-\thinmuskip
  \int#1%
  \ifx\\#3\\\else_{#3}\fi
  \ifx\\#4\\\else^{#4}\fi  
}
\newcommand*{\mint@@}[7]{%
  % #1: limits
  % #2: overlay symbol
  % #3: subscript
  % #4: superscript
  % #5: math style
  % #6: math style for overlay symbol
  % #7: math style for subscript/superscript
  \begingroup
    \sbox0{$#5\int\m@th$}%
    \sbox2{$#5\int_{}\m@th$}%
    \dimen2=\wd0 %
    % => \dimen2 = width of \int
    \let\mint@limits=#1\relax
    \ifx\mint@limits\relax
      \sbox4{$#5\int_{\kern1sp}^{\kern1sp}\m@th$}%
      \ifdim\wd4>\wd2 %
        \let\mint@limits=\nolimits
      \else
        \let\mint@limits=\limits
      \fi
    \fi
    \ifx\mint@limits\displaylimits
      \ifx#5\displaystyle
        \let\mint@limits=\limits
      \fi
    \fi
    \ifx\mint@limits\limits
      \sbox0{$#7#3\m@th$}%
      \sbox2{$#7#4\m@th$}%
      \ifdim\wd0>\dimen2 %
        \dimen2=\wd0 %
      \fi
      \ifdim\wd2>\dimen2 %
        \dimen2=\wd2 %
      \fi
    \fi
    \rlap{%
      $#5%
        \vcenter{%
          \hbox to\dimen2{%
            \hss
            $#6{#2}\m@th$%
            \hss
          }%
        }%
      $%
    }%
  \endgroup
}
\newcommand{\avg}{\mint{-}}
\makeatother

% ============================================================
% Comments and drafting helpers
% ============================================================

%Comments
\newcommand{\com}[1]{\vspace{5 mm}\par \noindent
\marginpar{\textsc{Comment}} \framebox{\begin{minipage}[c]{0.95
\textwidth} \tt #1 \end{minipage}}\vspace{5 mm}\par}

% ============================================================
% Arrows
% ============================================================

\newcommand{\hooktwoheadrightarrow}{%
  \hookrightarrow\mathrel{\mspace{-15mu}}\rightarrow
}

\makeatletter
\newcommand{\xhooktwoheadrightarrow}[2][]{%
  \lhook\joinrel
  \ext@arrow 0359\rightarrowfill@ {#1}{#2}%
  \mathrel{\mspace{-15mu}}\rightarrow
}
\makeatother

\makeatletter
\newcommand*{\da@rightarrow}{\mathchar"0\hexnumber@\symAMSa 4B }
\newcommand*{\da@leftarrow}{\mathchar"0\hexnumber@\symAMSa 4C }
\newcommand*{\xdashrightarrow}[2][]{%
  \mathrel{%
    \mathpalette{\da@xarrow{#1}{#2}{}\da@rightarrow{\,}{}}{}%
  }%
}
\newcommand{\xdashleftarrow}[2][]{%
  \mathrel{%
    \mathpalette{\da@xarrow{#1}{#2}\da@leftarrow{}{}{\,}}{}%
  }%
}
\newcommand*{\da@xarrow}[7]{%
  % #1: below
  % #2: above
  % #3: arrow left
  % #4: arrow right
  % #5: space left 
  % #6: space right
  % #7: math style 
  \sbox0{$\ifx#7\scriptstyle\scriptscriptstyle\else\scriptstyle\fi#5#1#6\m@th$}%
  \sbox2{$\ifx#7\scriptstyle\scriptscriptstyle\else\scriptstyle\fi#5#2#6\m@th$}%
  \sbox4{$#7\dabar@\m@th$}%
  \dimen@=\wd0 %
  \ifdim\wd2 >\dimen@
    \dimen@=\wd2 %   
  \fi
  \count@=2 %
  \def\da@bars{\dabar@\dabar@}%
  \@whiledim\count@\wd4<\dimen@\do{%
    \advance\count@\@ne
    \expandafter\def\expandafter\da@bars\expandafter{%
      \da@bars
      \dabar@ 
    }%
  }%  
  \mathrel{#3}%
  \mathrel{%   
    \mathop{\da@bars}\limits
    \ifx\\#1\\%
    \else
      _{\copy0}%
    \fi
    \ifx\\#2\\%
    \else
      ^{\copy2}%
    \fi
  }%   
  \mathrel{#4}%
}
\makeatother

% ============================================================
% Relation and delimiter spacing
% ============================================================

\renewcommand{\geq}{\geqslant}
\renewcommand{\leq}{\leqslant}
\newcommand{\midd}{\hspace{4pt}\middle|\hspace{4pt}}

% ============================================================
% Left Right Updated Deliminators
% ============================================================

\makeatletter
\NewDocumentCommand\Left{O{}}{%
    \IfValueTF{#1}{\notblank{#1}{\mathopen\bBigg@{#1}}{\left}}{\left}}
\NewDocumentCommand\Right{O{}}{%
    \IfValueTF{#1}{\notblank{#1}{\mathopen\bBigg@{#1}}{\right}}{\right}}
\makeatother

% ============================================================
% Restrictions
% ============================================================

\newcommand{\littletaller}{\mathchoice{\vphantom{\big|}}{}{}{}}
\newcommand\restr[2]{{% we make the whole thing an ordinary symbol
  \left.\kern-\nulldelimiterspace % automatically resize the bar with \right
  #1 % the function
  \littletaller % pretend it's a little taller at normal size
  \right|_{#2} % this is the delimiter
  }}

% ============================================================
% Chapter title formatting
% ============================================================

\newcommand{\chnum}{\titleformat
{\chapter} % command
[display] % shape
{\centering} % format
{\Huge \color{black} \shadowbox{\thechapter}} % label
{-0.5em} % sep (space between the number and title)
{\LARGE \color{black} \underline} % before-code
}

\newcommand{\chunnum}{\titleformat
{\chapter} % command
[display] % shape
{} % format
{} % label
{0em} % sep
{ \begin{flushright} \begin{tabular}{r}  \Huge \color{black}
} % before-code
[
\end{tabular} \end{flushright} \normalsize
] % after-code
}


\begin{document}
\homeworktitle
{Math 5023} %class
{Homework 1} %assignment
{Gianluca Crescenzo} %name

%%%%%%%%%%%%%%%%%%%%%%%%%%%%%%%%%%%
%%%%%%%%%%%%%%%%%%%%%%%%%%%%%%%%%%%
%%%%%%%%%%%%%%%%%%%%%%%%%%%%%%%%%%%
Throughout, let $R$ be a commutative ring with $1$, and $M$ an $R$-module.

\begin{exercise}
	Let $I \subset R$ be an ideal, $M$ an $R$-module. Prove $IM = \left\{\sum a_i m_i : a_i \in I, m_i \in M,\text{ finite sum}\right\}$ is a submodule of $M$.
\end{exercise}
{\color{blue}

	\begin{proof}
		It is clear that $IM$ is nonempty and closed under addition, as $1 \in IM$ and the sum of two finite sums will also be finite and of the form $\sum_{}^{} a_i m_i$. For any $r \in R$ and $\sum_{}^{} a_i m_i$, we have $r \left( \sum_{}^{} a_i m_i \right) = \sum_{}^{} (ra_i)m_i$. Since $I$ is an ideal, $ra_i \in I$. Hence $IM$ is a submodule of $M$
	\end{proof}
}
\begin{exercise}
	Suppose $M$ is an $R$-module, and $X \subset M$ is a subset. Prove that
$N = \langle X\rangle$ is the \textit{smallest} submodule containing the set $X$.
First formulate what \lq\lq smallest\rq\rq\ means.
\end{exercise}
{\color{blue}

	\begin{proof}
		Suppose $Y$ is a submodule of $M$ such that $X \subset Y \subset N$. Let $\sum_{}^{} rx \in N$. Since $X \subseteq Y$, $x \in Y$. Since $Y$ is a submodule, $\sum_{}^{} rx \in Y$. Thus $N \subseteq Y$; i.e., $Y = N$.
	\end{proof}
}
\begin{exercise}
	Let $V$ be the 3-dimensional $\mathbb{Q}$-vector space $\{a_0+a_1\sqrt[3]{2}+a_2\sqrt[3]{4}:a_i\in\mathbb{Q}\}$, and let $T=\sqrt[3]{2}+\sqrt[3]{4}$.
	\mbox{}
	\begin{enumerate}[label = (\alph*),itemsep=1pt,topsep=3pt]
		\item Convince yourself that left-multiplication by $T$ is a $\bfQ$-linear transformation on $V$. 
		\item Find a matrix for $T$ with respect to the basis $\{1,\sqrt[3]{2},\sqrt[3]{4}\}$.
		\item Compute $f(x)v$, where $f(x)=x^2-2x-4$, $x$ acts as the linear transformation $T$, and $v=2+\sqrt[3]{2}+\sqrt[3]{4}$.
		\item Find a polynomial $p(x)$ such that $p(x)v=0$ for all $v$.
	\end{enumerate}
\end{exercise}
{\color{blue}

	\begin{proof}
		(a) Given \(T=\sqrt[3]{2}+\sqrt[3]{4}\) and \(v\in V\), write, \(v=a_0+a_1\sqrt[3]{2}+a_2\sqrt[3]{4}\), where \(a_0,a_1,a_2\in\bfQ\). Then:
		\begin{align*}
			T(v)
&= \left(\sqrt[3]{2}+\sqrt[3]{4}\right)
   \left(a_0+a_1\sqrt[3]{2}+a_2\sqrt[3]{4}\right) \\
&= a_0\sqrt[3]{2}+a_0\sqrt[3]{4}
   +a_1\sqrt[3]{4}+2a_1+2a_2+2a_2\sqrt[3]{2} \\
&= (2a_1+2a_2)+(a_0+2a_2)\sqrt[3]{2}
   +(a_0+a_1)\sqrt[3]{4} \\
&\in V.
		\end{align*}
		Since each coefficient is rational, \(T\) takes elements of \(V\) to elements of \(V\). Moreover, by associativity and distributivity, for any \(v_1,v_2\in V\) and \(q\in\mathbb{Q}\), we have $T(v_1+qv_2)=Tv_1+qTv_2$. Therefore, left-multiplication by \(T\) is a \(\bfQ\)-linear transformation on \(V\).

		(b) Let \(\mathcal{B}=(1,\sqrt[3]{2},\sqrt[3]{4})\). We compute the image of each basis vector as follows:
		\begin{align*}
			T(1)
&= 0\cdot 1+1\cdot\sqrt[3]{2}+1\cdot\sqrt[3]{4}, \\
T\left(\sqrt[3]{2}\right)
&= 2\cdot 1+0\cdot\sqrt[3]{2}+1\cdot\sqrt[3]{4}, \\
T\left(\sqrt[3]{4}\right)
&= 2\cdot 1+2\cdot\sqrt[3]{2}+0\cdot\sqrt[3]{4}.
		\end{align*}
		Thus:
		\[
			[T]_{\mathcal{B}}
=
\begin{pmatrix}
0 & 2 & 2 \\
1 & 0 & 2 \\
1 & 1 & 0
\end{pmatrix}.
		\] 

		(c) We have:
		\begin{align*}
			f(x)v
&= (T^2-2T-4)\left(2+\sqrt[3]{2}+\sqrt[3]{4}\right) \\
&= \left(\left(\sqrt[3]{2}+\sqrt[3]{4}\right)^2
-2\left(\sqrt[3]{2}+\sqrt[3]{4}\right)-4\right)
\left(2+\sqrt[3]{2}+\sqrt[3]{4}\right) \\
&= \left(\sqrt[3]{4}+4+2\sqrt[3]{2}
-2\sqrt[3]{2}-2\sqrt[3]{4}-4\right)
\left(2+\sqrt[3]{2}+\sqrt[3]{4}\right) \\
&= -\sqrt[3]{4}\left(2+\sqrt[3]{2}+\sqrt[3]{4}\right) \\
&= -2\sqrt[3]{4}-2-2\sqrt[3]{2}.
		\end{align*}

		(d) Let \(p(x)\) be the characteristic polynomial of \([T]_{\mathcal{B}}\). Then \(p(x)v=0\) for all \(v\in V\) by the Cayley--Hamilton theorem.
	\end{proof}
}

\begin{exercise}
	Suppose $\{M_i:i\in I\}$ is a set of submodules of $M$
	\mbox{}
	\begin{enumerate}[label = (\alph*),itemsep=1pt,topsep=3pt]
		\item Prove that $\sum_i M_i$ is a submodule of $M$, where all sums are finite. Prove that $\sum_i M_i$ is the smallest submodule of $M$ that contains all of the $M_i$.:w
		\item Prove that $\bigcap_i M_i$ is a submodule of $M$.
	\end{enumerate}
\end{exercise}
{\color{blue}

	\begin{proof}
		Define $N := \sum_{i}^{} M_i$.

		(a) Note $N \neq 0$ since $1 \in N$. Let $\sum_{}^{} m_i, \sum_{}^{n_i} \in N$. The sum of these two elements is clearly a finite sum, hence they are an element of $N$. Let $r \in R$ and $\sum_{}^{} m_i \in N$. Then $r \left( \sum_{}^{} m_i \right)  = \sum_{}^{} r m_i$. Since $rm_i \in M_i$ for each $i$, we have $r\left( \sum_{}^{} m_i \right) \in N$. Thus $N$ is a submodule of $M$. Suppose now that $Y$ is a submodule of $M$ such that $\bigcup_{i} M_i \subseteq Y \subseteq N.$ Let $\sum_{}^{}m_i \in N$. Note that each $m_i \in \bigcup_i M_i \subseteq Y$. Since $Y$ is a submodule, $\sum_{}^{} m_i \in Y$. Hence $Y = N$.

		(b) Note $\bigcap_i M_i \neq 0 $ because $1 \in M_i$ for all $i$. Let $m,n \in \bigcap_i M_i$. Then $m,n \in M_i$ for all $i$. Hence for all $r \in R$, we have $m+rn \in M_i$ for all $i$. Thus $m+rn \in \bigcap_i M_i$, hence it is a submodule of $M$.
	\end{proof}
}

\begin{exercise}
	Suppose $M_1\subset M_2\subset M_3\subset\cdots$ is an ascending sequence of submodules of an $R$-module $M$. Prove that $\bigcup_i M_i$ is a submodule of $M$.
\end{exercise}
{\color{blue}

	\begin{proof}
		Note that $1 \in \bigcup_{i}M_i$ because $1 \in M_i$ for all $i$. If $m,n \in \bigcup_i M_i$, then $m \in M_j$ and $n \in M_k$ for some indices $j,k$. There exists some index $\ell$ such that $M_j,M_k \subseteq M_\ell$, hence $m,n \in M_\ell$. Thus $m+n \in M_\ell \subseteq \bigcup_i M_i$. Moreover, for any $r \in R$, we have $mr \in M_\ell \subseteq \bigcup_i M_i$. Thus $\bigcup_i M_i$ is a submodule.
	\end{proof}
}

\begin{exercise}
	Suppose $f : M \longrightarrow M'$ and $g : N \longrightarrow N'$ are
$R$-module homomorphisms.
	\mbox{}
	\begin{enumerate}[label = (\alph*),itemsep=1pt,topsep=3pt]
		\item Prove that $f$ and $g$ induce $R$-module homomorphisms $f^* : \mathrm{Hom}_R(M',N) \longrightarrow \mathrm{Hom}_R(M,N)$ and $g_* : \mathrm{Hom}_R(M,N) \longrightarrow \mathrm{Hom}_R(M,N')$.
		\item Prove that if $f$ is onto then $f^*$ is one-to-one, and if $g$ is one-to-one then $g_*$ is one-to-one.
	\end{enumerate}
\end{exercise}
{\color{blue}

	\begin{proof}
		(a) Define $f^\ast:\Hom_R(M',N) \rightarrow \Hom_R(M,N)$ by $f^\ast(\varphi) := f \circ \varphi$. This is indeed an $R$-module homomorphism, as for any $\varphi_1,\varphi_2 \in \Hom_R(M',N)$ and $r \in R$ we have:
		\begin{align*}
			f^\ast(\varphi_1 + r\varphi_2)(m) 
			&= \bigl( (\varphi_1 + r \varphi_2) \circ f \bigr)(m) \\
			&=  (\varphi_1 + r \varphi_2)\bigl( f(m) \bigr)\\
			&= \varphi_1(f(m)) + r \varphi_2(f(m)) \\
			&= \bigl( (\varphi_1 \circ f) + r (\varphi_2 \circ f) \bigr)(m) \\
			&= \bigl( f^\ast(\varphi_1) + r f^\ast(\varphi_2) \bigr)(m).
		\end{align*}
		Similarly, define $g_\ast:\Hom_R(M,N) \rightarrow \Hom_R(M,N')$ by $g_\ast(\varphi):= g \circ \varphi$. Then $g_\ast$ is an $R$-module homomorphism, as for any $\varphi_1,\varphi_2 \in \Hom_R(M,N)$ and $r \in R$ we have:
		\begin{align*}
			g_\ast(\varphi_1 + r \varphi_2)(m)
			&= \bigl( g \circ (\varphi_1 + r \varphi_2) \bigr)(m) \\
			&= g \bigl( (\varphi_1 + r\varphi_2)(m) \bigr) \\
			&= g \bigl( \varphi_1(m) + r \varphi_2(m) \bigr) \\
			&= (g \circ \varphi_1)(m) + r (g \circ \varphi_2)(m) \\
			&= \bigl( (g \circ \varphi_1 \bigr) + r (g \circ \varphi_2)(m).
		\end{align*}

		(b) Suppose $f^\ast(\varphi_1) = f^\ast(\varphi_2)$ for some $\varphi_1,\varphi_2 \in \Hom_R(M',N)$. Then $\varphi_1(f(m)) = \varphi_2(f(m))$ for every $m \in M'$. Since $f$ is onto, there exists $x \in M$ such that $f(m) = x$. Thus $\varphi_1(x) = \varphi_2(x)$. Since $\varphi_1$ and $\varphi_2$ are equal on all elements, we have that $\varphi_1 = \varphi_2$. Thus $f^\ast$ is one-to-one. Now suppose $g_\ast(\psi_1) = g_\ast(\psi_2)$ for some $\psi_1,\psi_2 \in \Hom_R(M,N)$. Then $g(\psi_1(m)) = g(\psi_2(m))$ for all $m \in M$. Since $g$ is one-to-one, $\psi_1(m) = \psi_2(m)$ for all $m \in M$. Since $\psi_1$ and $\psi_2$ are equal on all elements, we have $\psi_1 = \psi_2$. Thus $g_\ast$ is one-to-one.
	\end{proof}
}

\begin{exercise}
	If $N$ is a submodule of $M$, the \textit{annihilator of $N$ in $R$} is
$\mathrm{Ann}_R(N)=\{r\in R:rn=0\ \forall n\in N\}$.
	\mbox{}
	\begin{enumerate}[label = (\alph*),itemsep=1pt,topsep=3pt]
		\item Prove $\mathrm{Ann}_R(N)$ is a 2-sided ideal of $R$. 
		\item Suppose $M=\mathbb{Z}_{24}\times\mathbb{Z}_{15}\times\mathbb{Z}_{50}$. Compute $\mathrm{Ann}_{\mathbb{Z}}(M)$.
	\end{enumerate}
\end{exercise}
{\color{blue}

	\begin{proof}
		(a) Note that $\Ann_R(N) \neq \emptyset$ because $0 \in \Ann_R(N)$. If $r_1,r_2 \in \Ann_R(N)$, then $(r_1 - r_2)n = r_1n - r_2n = 0$, hence $r_1 - r_2 \in Ann_R(N)$. If $x \in \Ann_R(N)$, and $r \in R$, then $(xr)n = (rx)n = 0$ and $(rx)n =  r(xn) = 0$. Thus $xr,rx \in \Ann_R(N)$, implying that it is a 2-sided ideal of $R$;

		(b) It is enough to determine all $\alpha \in \bfZ$ such that $\alpha(\overline{1},\overline{0},\overline{0}) = 0$, $\alpha(\overline{0},\overline{1},\overline{0}) = 0$, and $\alpha(\overline{0},\overline{0},\overline{1}) = 0$. This is equivalent to determining all $\alpha$ such that $\overline{\alpha} \equiv 0 \pmod{24}$, $\overline{\alpha}\equiv 0\pmod{15}$, and $\overline{\alpha}\equiv 0 \pmod{50}$. This system of equations is equivalent to $\overline{\alpha} \equiv 0 \pmod{\lcm(24,15,50)}$, thus $\Ann_\bfZ(M) = \left<\lcm(24,15,50) \right>$.
	\end{proof}
}

\begin{exercise}
	Assume $R$ is commutative. Show that $M$ is an irreducible $R$-module if and only if $M$ is $R$-isomorphic to $R/\mathfrak m$, for some maximal (2-sided) ideal $\mathfrak m$.
\end{exercise}
{\color{blue}

	\begin{proof}
		Suppose that $M$ is an irreducible $R$-module. For any $m \in M$, note that $\left< m\right>$ is a submodule of $M$. Hence $M = \left<m \right>$. Define $\varphi:R \rightarrow M$ by $r \mapsto rm$. This is clearly an $R$-module homomorphism, as $\varphi(r_1 + r_2) = (r_1 + r_2)m = r_1m + r_2m = \varphi(r_1) + \varphi(r_2)$ and $\varphi(r r_0) = (r r_0)m = r (r_0 m) = r \varphi(r_0).$ Treating $M$ and $R$ as groups, by the First Isomorphism Theorem we have $R/\ker\varphi \cong \Image \varphi$. Note that $\varphi$ is surjective, as for any $m \in M = \left<m \right>$, we have $m = r_0m$ for some $r_0 \in R$; i.e., $\varphi(r_0) = r_0m = m$. Hence $R/\ker\varphi = M$. It remains to show that $\ker \varphi$ is a maximal ideal. Suppose $J$ is an ideal of $R$ such that $\ker\varphi \subseteq J \subseteq R$. By the Third Isomorphism Theorem, we have $(R/\ker\varphi) (J/\ker\varphi) \cong R/J$, and in particular, $J/\ker\varphi$ is an ideal of $R/\ker\varphi \cong M$. Since $M$ is irreducible, either $J/\ker\varphi = 0$ or $J/\ker\varphi = M$. If $J/\ker\varphi = 0$, then $J = \ker\varphi$. If $J/\ker\varphi = M$, then $J = R$. Thus $\ker\varphi$ is a maximal ideal.

		Conversely, let $\varphi: M \rightarrow R/\fm$ be an $R$-isomorphism for some maximal ideal $\fm$. Let $N$ be a submodule of $M$. Note that $\varphi(N)$ is an ideal, as it is closed under addition and scalar multiplication by $R$. Since $\fm$ is maximal, $R/\fm$ is a field. Hence $\varphi(N)=  0$ or $\varphi(N) = R/\fm$, implying that $N = 0$ or $N = M$. Thus $M$ is irreducible.
	\end{proof}
}

\begin{exercise}
	Suppose $N\subset M$ is a submodule. Prove that if $N$ and $M/N$ are finitely generated, then $M$ is finitely generated.
\end{exercise}
{\color{blue}

	\begin{proof}
		If $N$ is finitely generated, then $N = \left<X \right>$ for some $X \subset N$. Let $m \in M$ be arbitrary and consider $m +\left<X \right> \in M / \left<X \right>$. Since $M / \left<X \right>$ is finitely generated, there is a subset $Y \subset M/\left<X \right>$ such that $M/\left<X \right> = \left<Y \right>$. Hence there exists $r_1,\ldots ,r_n \in R$ and $y_1+\left<X \right>,\ldots ,y_n+\left<X \right> \in Y$ such that:
		\begin{align*}
			m + \left<X \right>
			&= \sum_{i=1}^{n} r_i (y_i + \left<X \right>) \\
			&= \left(\sum_{i=1}^{n} r_i y_i \right) + \left<X \right>.
		\end{align*}
		This means $m - \left( \sum_{i=1}^{n} r_i y_i \right) \in \left<X \right>$, so there exists $s_1,\ldots ,s_k \in R$ and $x_1,\ldots ,x_k \in X$ such that:
		\[
			m - \left( \sum_{i=1}^{n} r_i y_i \right) = \sum_{i=1}^{k} s_i x_i.
		\] 
		Thus $m = \left( \sum_{i=1}^{n} r_i y_i \right) + \left( \sum_{i=1}^{k} s_i x_i \right)$. Since $m \in M$ was arbitary, $M$ is finitely generated.
	\end{proof}
}



\end{document}
